% Generated draft: requires mathematical and source review.
\documentclass[11pt]{article}
\usepackage[margin=25mm]{geometry}
\usepackage{fontspec}
\setmainfont{Latin Modern Roman}
\usepackage{amsmath,amssymb,mathtools,unicode-math}
\setmathfont{Latin Modern Math}
\usepackage{enumitem,needspace,xurl,hyperref}
\setlength{\parindent}{0pt}
\setlength{\parskip}{5pt}
\setlength{\emergencystretch}{3em}
\newcommand{\N}{\mathbb{N}}
\newcommand{\Z}{\mathbb{Z}}
\newcommand{\Q}{\mathbb{Q}}
\newcommand{\R}{\mathbb{R}}
\newcommand{\C}{\mathbb{C}}
\newcommand{\F}{\mathbb{F}}
\newcommand{\A}{\mathbb{A}}
\newcommand{\PP}{\mathbb{P}}
\newcommand{\E}{\mathbb{E}}
\newcommand{\eps}{\varepsilon}
\newcommand{\dd}{\,\mathrm{d}}
\begin{document}
\section*{UnsolvedMath ID: 30004769}
\textbf{Source problem code:} catalogue entry
\subsection{Definitions and mathematical statement}
For a group \(G\) acting on a finite set \(X\), the orbit of \(x\in X\) is
\[
G\cdot x=\{g\cdot x:g\in G\},
\]
and the orbit set is \(G\backslash X=\{G\cdot x:x\in X\}\). Determining the orbits means determining this partition of \(X\), for example by giving one representative of each orbit and a procedure assigning each element to its orbit.

The supplied record concerns an infinite group \(G\) and finite quotients denoted by
\[
Q_{n,e}=\frac{(\theta-1)^n\Gamma}{(\theta-1)^{n+e}\Gamma},
\]
arising from a construction of constructible \(p\)-groups of maximal class. Here \(p\) denotes a prime; the specific construction data and admissible parameter values must be supplied by that construction. For context, a finite group of order \(p^m\), with \(m\geq2\), has maximal class when its nilpotency class is \(m-1\). This does not define the source-specific term ``constructible.''

\textbf{Computational problem.} For each admissible pair \((n,e)\) in the intended construction, develop computational methods to determine
\[
G\backslash Q_{n,e}
\]
for the specified action of \(G\) on \(Q_{n,e}\), allowing \(n\) and \(e\) to vary.

\textbf{Unresolved qualifications.} The supplied record does not define \(\Gamma\), \(\theta\), the infinite group \(G\), or its action. It does not specify the admissible values of \(p,n,e\), the meaning of ``constructible'' in this construction, or effective representations of the objects and action. Thus the displayed quotient and its action are retained as source notation, without asserting additional algebraic structure or parameter bounds. The problem above is a conditional formulation of the stated computational target; these missing data prevent a fully specified algorithmic problem. No complexity bound or particular form of computational method is supplied.
\subsection{Short English statement}
Develop computational methods to find the orbits of the specified infinite group on the finite quotients \((\theta-1)^n\Gamma/(\theta-1)^{n+e}\Gamma\) as the admissible parameters vary. The construction, action, parameter restrictions, and computational input model are missing from the supplied record and require clarification from the cited report.
\subsection{Sources}
\begin{itemize}
\item Local registry: \texttt{data/unsolvedmath/problems.json}, permanent UnsolvedMath ID \(30004769\), title ``Orbit Computation for Constructible Maximal-Class p-Groups,'' source problem code \texttt{catalogue entry}; collection: Oberwolfach Reports Open Problems. No individual proposer is identified; the supplied proposed year is 2021.
\item Supplied report citation: \emph{Computational Group Theory (2022)}, Oberwolfach Reports, pp.\ 2027--2087, DOI \texttt{10.4171/owr/2021/38}. Supplied source URL: \url{https://doi.org/10.4171/owr/2021/38}. The literature-source entry identifies the same link as ``Oberwolfach Report 38/2021.'' Both the supplied 2022 citation designation and the 2021 report designation are retained without resolving their bibliographic relationship.
\item Supplied literature assessment, checked \textbf{2026-08-22}: the record labels the problem open, describes the report as presenting computational investigation of the orbit problem as an open direction, and says that no general algorithmic resolution was found. Its listed literature source is \url{https://doi.org/10.4171/owr/2021/38}. This is supplied dated curation, not an independently verified literature assessment.
\item The registry labels the statement \texttt{corrected\_verified} and reports that the official report was checked directly and the clean statement restored definitions while removing extraction artifacts or source-dependent wording. These are supplied verification claims; this draft does not independently verify them, and the supplied text does not actually include the construction definitions needed above.
\end{itemize}
\end{document}
